% answers_short.tex —— 1979 年·全国卷（理科） · 简短答案速查
% 由 tools/make_short_answers.py 自 answers.tex 的 \gksolhead/\gkanswer 提取，
% 并与 1979/1979_数学真题_理_答案要点.md 逐题比对后生成，勿手改（改源 MD 后重跑）。

\begin{gknotice}
\textbf{使用说明}：本节只给最终答案便于快速核对，完整推导见解析版。答案为编写者独立推导，\textbf{非官方标准答案}。
\end{gknotice}

\gkpart{题目}
\begin{itemize}[leftmargin=2.6em]
  \item[\textbf{1.}] 见解析版
  \item[\textbf{2.}] $\csc^2x$
  \item[\textbf{3.}] $V_1m_1(m_2+n_2)+V_2m_2(m_1+n_1)\,:\,V_1n_1(m_2+n_2)+V_2n_2(m_1+n_1)$
  \item[\textbf{4.}] 见解析版
  \item[\textbf{5.}] 当 $\dfrac{S\sin\alpha\sin\beta}{\sin(\alpha+\beta)}<D$ 时发出警告；等价地 $\dfrac{\sin\alpha\sin\beta}{\sin(\alpha+\beta)}<\dfrac DS$
  \item[\textbf{6.}] 见解析版
  \item[\textbf{7.}] $x\approx0.04025$，即每年约增长 $4\%$
  \item[\textbf{8.}] 见解析版
  \item[\textbf{9.}] 前 $14$ 项的和最大，最大值为 $14.3045$
  \item[\textbf{10.}] (1) 轨迹为圆 $\left(x-h\sec^2\theta\right)^2+y^2=\left(h\tan\theta\sec\theta\right)^2$ 位于 $\triangle OAB$ 内的弧 $AB$（含端点 $A$、$B$）；(2) $P\left(\dfrac{h\sqrt5}{\sqrt5+2\sin\theta},\ \dfrac{h\tan\theta}{\sqrt5+2\sin\theta}\right)$
\end{itemize}
